\documentclass[11pt]{amsart}

\usepackage[T1]{fontenc}
\usepackage{lmodern}
\usepackage{microtype}
\usepackage{amsmath,amssymb,amsthm}
\usepackage[hidelinks]{hyperref}

\hypersetup{
  pdftitle={A Counterexample to the Gallian--Schoenhard Star-Forest Conjecture}
}

\newtheorem{theorem}{Theorem}
\newtheorem{lemma}[theorem]{Lemma}
\newtheorem{corollary}[theorem]{Corollary}
\theoremstyle{remark}
\newtheorem{remark}[theorem]{Remark}

\title[A counterexample to the Gallian--Schoenhard conjecture]
{A Counterexample to the Gallian--Schoenhard Star-Forest Conjecture}

\date{}

\subjclass[2020]{05C78}
\keywords{
even harmonious labeling,
strongly even harmonious labeling,
star forest
}

\begin{document}

\begin{abstract}
Gallian and Schoenhard conjectured that a disjoint union of stars is
strongly even harmonious whenever one component has more than two
edges. We disprove the conjecture: $S_3\sqcup 3S_1$ is properly even
harmonious but not strongly even harmonious. The nonexistence holds
even with vertex labels in the larger codomain $\{0,\ldots,2q\}$, and
therefore also under the corrected proper codomain
$\{0,\ldots,2q-1\}$. A counting inequality also gives two infinite
families of counterexamples.
\end{abstract}

\maketitle

\section{The counterexample}

Write $S_n=K_{1,n}$, so $S_n$ has $n$ edges, and let $\sqcup$ denote
disjoint union. In the original even-harmonious convention, a graph
$G$ with $q$ edges is strongly even harmonious if there is an
injection
\[
 f:V(G)\longrightarrow \{0,1,\ldots,2q\}
\]
such that the induced edge labels
\[
 f^*(uv)=f(u)+f(v)\pmod{2q}
\]
are exactly $0,2,\ldots,2q-2$, and every edge $uv$ satisfies
\[
 0<f(u)+f(v)\le 2q.
\]
This is the convention used in the original paper and retained by
Henderson~\cite{GallianSchoenhard,Henderson}. Gallian's current survey
instead builds the strong condition on the corrected properly even
harmonious codomain
$\{0,1,\ldots,2q-1\}$~\cite[Section~4.6]{Gallian}; see also
Henderson~\cite{Henderson}. Since this is the smaller codomain, a
nonexistence result for $\{0,\ldots,2q\}$ applies under both
conventions.

Gallian and Schoenhard conjectured that
\[
 S_{n_1}\sqcup\cdots\sqcup S_{n_t}
\]
is strongly even harmonious whenever at least one $n_i$ exceeds
$2$~\cite[Conjecture~4.12]{GallianSchoenhard}. The statement is still
listed as a conjecture in the 2025 edition of Gallian's
survey~\cite[Section~4.6, p.~142]{Gallian}.

\begin{theorem}\label{thm:main}
The star forest
\[
 G=S_3\sqcup 3S_1
\]
is properly even harmonious but has no strongly even harmonious
labeling, even with vertex labels in $\{0,1,\ldots,12\}$. Consequently,
Conjecture~4.12 of Gallian and Schoenhard is false under either
codomain convention.
\end{theorem}

The proof uses the following counting obstruction.

\begin{lemma}\label{lem:count}
Let
\[
 F=S_{n_1}\sqcup\cdots\sqcup S_{n_t},
 \qquad
 n_1\ge\cdots\ge n_t\ge1,
 \qquad
 q=\sum_{i=1}^t n_i.
\]
If $F$ has a strongly even harmonious labeling with vertex labels in
$\{0,1,\ldots,2q\}$, then
\begin{equation}\label{eq:obstruction}
 q^2+\sum_{i=1}^t(i-1)(n_i-1)
 \le
 \frac{(q-t+1)(3q+t)}{2}.
\end{equation}
\end{lemma}

\begin{proof}
Let $c_i$ be the center of $S_{n_i}$, and let $U$ be the sum of the
unused vertex labels. The strong bound lifts the edge residues
$0,2,\ldots,2q-2$ uniquely to the ordinary edge sums
\[
 2,4,\ldots,2q,
\]
whose total is $q(q+1)$. Therefore
\[
 q(q+1)
 =
 \sum_{v\in V(F)}f(v)
 +
 \sum_{i=1}^t(n_i-1)f(c_i).
\]
Since the sum of all labels in $\{0,\ldots,2q\}$ is $q(2q+1)$,
\begin{equation}\label{eq:unused}
 U=q^2+\sum_{i=1}^t(n_i-1)f(c_i).
\end{equation}

Put $w_i=n_i-1$, so $w_1\ge\cdots\ge w_t\ge0$, and let
$b_1<\cdots<b_t$ be the center labels in increasing order. Then
$b_i\ge i-1$, and the rearrangement inequality gives
\[
 \sum_{i=1}^t(n_i-1)f(c_i)
 \ge
 \sum_{i=1}^t w_i b_i
 \ge
 \sum_{i=1}^t(i-1)(n_i-1).
\]
There are $q-t+1$ unused labels. Their sum is at most the sum of the
$q-t+1$ largest labels in the codomain, namely
\[
 U
 \le
 (q+t)+(q+t+1)+\cdots+2q
 =
 \frac{(q-t+1)(3q+t)}{2}.
\]
Combining this bound with \eqref{eq:unused} proves
\eqref{eq:obstruction}.
\end{proof}

\begin{proof}[Proof of Theorem~\ref{thm:main}]
For $G=S_3\sqcup3S_1$, we have $q=6$, $t=4$, and
\[
 (n_1,n_2,n_3,n_4)=(3,1,1,1).
\]
Thus the two sides of \eqref{eq:obstruction} are $36$ and $33$,
respectively. Hence no strongly even harmonious labeling exists, even
in the larger codomain $\{0,\ldots,12\}$.

It remains to exhibit a properly even harmonious labeling. Let $c$ be
the center of $S_3$, with leaves $x_1,x_2,x_3$, and let $a_i b_i$,
$1\le i\le3$, be the three $S_1$ components. Set
\[
 f(c)=0,
 \qquad
 \bigl(f(x_1),f(x_2),f(x_3)\bigr)=(2,6,10),
\]
and
\[
 \bigl(f(a_1),f(b_1)\bigr)=(3,5),\qquad
 \bigl(f(a_2),f(b_2)\bigr)=(4,8),\qquad
 \bigl(f(a_3),f(b_3)\bigr)=(7,9).
\]
These are distinct labels in $\{0,\ldots,11\}$, and the induced edge
labels modulo $12$ are exactly
\[
 0,2,4,6,8,10.
\]
Thus $G$ is properly even harmonious.
\end{proof}

\begin{corollary}\label{cor:families}
Neither of the following star forests is strongly even harmonious
under either codomain convention:
\[
 S_3\sqcup mS_1 \quad (m\ge3),
 \qquad
 S_3\sqcup kS_2 \quad (k\ge5).
\]
\end{corollary}

\begin{proof}
For $S_3\sqcup mS_1$, the left side of
\eqref{eq:obstruction} exceeds the right side by
\[
 m^2-6>0.
\]
For $S_3\sqcup kS_2$, the weighted term in
\eqref{eq:obstruction} is $k(k+1)/2$, and the corresponding difference
is
\[
 k^2-3k-6>0.
\]
\end{proof}

\begin{remark}
Henderson proved that a union of at least two stars, each having at
least two edges, is properly even
harmonious~\cite[Theorem~3.1]{Henderson}. Consequently,
$S_3\sqcup kS_2$ is properly even harmonious but not strongly even
harmonious for every $k\ge5$.
\end{remark}

\begin{thebibliography}{9}

\bibitem{GallianSchoenhard}
J.~A. Gallian and L.~A. Schoenhard,
\emph{Even harmonious graphs},
AKCE Int. J. Graphs Comb. \textbf{11} (2014), no.~1, 27--49.
\href{https://doi.org/10.1080/09728600.2014.12088761}
{doi:10.1080/09728600.2014.12088761}.

\bibitem{Gallian}
J.~A. Gallian,
\emph{A dynamic survey of graph labeling},
Electron. J. Combin., Dynamic Survey \#DS6,
28th ed. (2025).
\href{https://doi.org/10.37236/27}
{doi:10.37236/27}.

\bibitem{Henderson}
Z.~M. Henderson,
\emph{Properly even harmonious labeling of a union of stars},
Electron. J. Graph Theory Appl. \textbf{12} (2024),
no.~1, 1--16.
\href{https://doi.org/10.5614/ejgta.2024.12.1.1}
{doi:10.5614/ejgta.2024.12.1.1}.

\end{thebibliography}

\end{document}